\section{Requisitos Arquitectónicamente Significativos (ASR)}
Los Requisitos Arquitectónicamente Significativos de Kit Pagos Colombia se derivaron cruzando los atributos de calidad que motivaron las decisiones de diseño ya documentadas (arquitectura hexagonal, normalización de respuestas, manejo de reintentos y verificación de firmas) con los requisitos funcionales definidos para el SDK y la API de Simulación. La tabla siguiente organiza estos requisitos en forma de árbol de utilidad, relacionando cada atributo de calidad con su sub-atributo específico, las funcionalidades que lo materializan y un escenario concreto que permite verificarlo. Los ASR están ordenados de mayor a menor prioridad, comenzando por aquellos que sustentan la propuesta central de la tesis (la abstracción unificada de múltiples pasarelas) y terminando por aquellos de carácter complementario.

Los ASR de prioridad alta se concentran en la extensibilidad ante nuevas pasarelas (ASR-01 y ASR-02) y la normalización uniforme de respuestas (ASR-03). También son de prioridad alta la autenticidad de los webhooks y la protección de credenciales (ASR-04 y ASR-07), y la reproducción fiel de escenarios junto con la recuperación ante fallas transitorias (ASR-05 y ASR-06). Los ASR de prioridad media y baja cubren el manejo uniforme de errores, el bajo acoplamiento, la disponibilidad del simulador y la documentación autoexplorable de la API.

Para que la tabla se pueda leer sin comprimir el texto ni reducir el tamaño de la fuente, los doce requisitos se adjuntan en un libro de Excel con formato, en \path{docs/final-doc/sad/anexos/asr.xlsx}. El encabezado va en negrita, la primera fila queda fija y el filtro permite ordenar por prioridad.

